% Part: first-order-logic
% Chapter: completeness
% Section: maximally-consistent-sets

% সর্বাধিক সঙ্গত সেটের সংজ্ঞা। পূর্ণতার জন্য প্রয়োজনীয় ধর্মগুলি ব্যবহৃত
% প্রমাণপদ্ধতির অধ্যায়ে আলোচনা করা হয়েছে।

\documentclass[../../../include/open-logic-section]{subfiles}

\begin{document}

\olfileid{fol}{com}{mcs}
\olsection{সর্বাধিক সঙ্গত \usetoken{P}{sentence}-সেট}

\begin{defn}[সর্বাধিক সঙ্গত সেট]
\ollabel{mcs}
!!{sentence}s-এর কোনো সেট~$\Gamma$ ঠিক তখনই \emph{সর্বাধিক সঙ্গত}, যখন
\begin{enumerate}
\item $\Gamma$ সঙ্গত, এবং
\item $\Gamma \subsetneq \Gamma'$ হলে $\Gamma'$ অসঙ্গত।
\end{enumerate}
\end{defn}

\begin{explain}
উপরের সংজ্ঞার সমতুল্য আরেকটি সংজ্ঞা হলো: বাক্যের কোনো সেট~$\Gamma$
ঠিক তখনই \emph{সর্বাধিক সঙ্গত}, যখন
\begin{enumerate}
\item $\Gamma$ সঙ্গত, এবং
\item $\Gamma \cup \{ !A \}$ সঙ্গত হলে $!A \in \Gamma$।
\end{enumerate}
অর্থাৎ, সর্বাধিক সঙ্গত সেট~$\Gamma$-তে আগে থেকে নেই এমন
!!{sentence}s সেই~$\Gamma$-তে যোগ করলে বড় সেটটি অসঙ্গত হয়ে যায়।
\end{explain}

\begin{explain}
পূর্ণতা-প্রমাণে সর্বাধিক সঙ্গত সেট গুরুত্বপূর্ণ। কারণ নিশ্চয়তা দেওয়া যায়,
!!{sentence}s-এর প্রত্যেক সঙ্গত সেট~$\Gamma$ কোনো সর্বাধিক সঙ্গত
সেট~$\Gamma^*$-এর অন্তর্ভুক্ত। আবার সর্বাধিক সঙ্গত সেটে প্রত্যেক
!!{sentence}~$!A$-এর জন্য $!A$ অথবা তার নঞর্থকরণ~$\lnot !A$ থাকে।
বিশেষত পরমাণু !!{sentence}s-এর ক্ষেত্রেও এটি সত্য। তাই যথাযথভাবে
!!{constant}s যোগ করে প্রসারিত ভাষায় কোনো সর্বাধিক সঙ্গত সেট থেকে
আমরা এমন একটি !!a{structure} নির্মাণ করতে পারি, যার
!!{predicate}s-এর ব্যাখ্যা নির্ধারিত হয়~$\Gamma^*$-এ কোন পরমাণু
!!{sentence}s আছে তার ভিত্তিতে। এরপর দেখানো যায়, ওই !!{structure}
$\Gamma^*$-এর সব !!{sentence}s-কে (এবং তাই $\Gamma$-র সবগুলিকেও)
সত্য করে। শেষের দাবিটি প্রমাণ করতে দরকার: $\lnot !A \in
\Gamma^*$ যদি এবং কেবল যদি $!A \notin \Gamma^*$; $(!A \lor !B) \in
\Gamma^*$ যদি এবং কেবল যদি $!A \in \Gamma^*$ অথবা $!B \in
\Gamma^*$; ইত্যাদি।
\end{explain}

\begin{prop}
\ollabel{prop:mcs}
ধরা যাক $\Gamma$ সর্বাধিক সঙ্গত। তাহলে:
% BN-SRC-860: unrequested restatement of the established sentence convention removed; source scope retained.
\begin{enumerate}
\item \ollabel{prop:mcs-prov-in} $\Gamma \Proves !A$ হলে $!A \in
  \Gamma$।

\item \ollabel{prop:mcs-either-or} যেকোনো $!A$-এর জন্য $!A \in
  \Gamma$ অথবা $\lnot !A \in \Gamma$।

\tagitem{prvAnd}{\ollabel{prop:mcs-and} $(!A \land !B) \in \Gamma$
  যদি এবং কেবল যদি $!A \in \Gamma$ এবং $!B \in \Gamma$ উভয়ই সত্য।}{}

\tagitem{prvOr}{\ollabel{prop:mcs-or} $(!A \lor !B) \in \Gamma$ যদি
  এবং কেবল যদি $!A \in \Gamma$ অথবা $!B \in \Gamma$।}{}

\tagitem{prvIf}{\ollabel{prop:mcs-if} $(!A \lif !B) \in \Gamma$ যদি
  এবং কেবল যদি $!A \notin \Gamma$ অথবা $!B \in \Gamma$।}{}
\end{enumerate}
\end{prop}

\begin{proof}
নিচের প্রতিটি ক্ষেত্রে ধরে নিই যে $\Gamma$ সর্বাধিক সঙ্গত।
\begin{enumerate}
\item $\Gamma \Proves !A$ হলে $!A \in \Gamma$।

ধরা যাক $\Gamma \Proves !A$। বিপরীতভাবে ধরি $!A
\notin \Gamma$। $\Gamma$ সর্বাধিক সঙ্গত বলে $\Gamma
\cup \{!A\}$ অসঙ্গত; অর্থাৎ $\Gamma \cup \{!A\} \Proves
\lfalse$।
\tagrefs{prfSC/{fol:seq:prv:prop:provability-contr},prfND/{fol:ntd:prv:prop:provability-contr}}
অনুসারে $\Gamma$ অসঙ্গত। এটি $\Gamma$-র সঙ্গতির পূর্বধারণার বিরোধী।
অতএব $!A \notin \Gamma$ হতে পারে না; তাই $!A \in \Gamma$।

\item যেকোনো $!A$-এর জন্য $!A \in \Gamma$ অথবা $\lnot !A \in \Gamma$।

বিপরীতভাবে ধরি, কোনো~$!A$-এর জন্য $!A \notin \Gamma$ এবং
$\lnot !A \notin \Gamma$ দুটিই সত্য। $\Gamma$ সর্বাধিক সঙ্গত বলে
$\Gamma \cup \{!A\}$ এবং $\Gamma \cup \{\lnot !A\}$ উভয়ই
অসঙ্গত। ফলে $\Gamma \cup \{!A\} \Proves \lfalse$ এবং $\Gamma \cup
\{\lnot !A\} \Proves \lfalse$।
\tagrefs{prfSC/{fol:seq:prv:prop:provability-exhaustive},prfND/{fol:ntd:prv:prop:provability-exhaustive}}
অনুসারে $\Gamma$ অসঙ্গত—এটি স্ববিরোধ। তাই এমন কোনো
!!a{sentence}~$!A$ হতে পারে না; প্রত্যেক $!A$-এর জন্য $!A \in
\Gamma$ অথবা $\lnot !A \in \Gamma$।

\tagitem{defAnd}{}{%
\iftag{probAnd}{অনুশীলনী।}{%
$(!A \land !B) \in \Gamma$ যদি এবং কেবল যদি $!A \in \Gamma$ ও
$!B \in \Gamma$ উভয়ই সত্য:

সামনের অভিমুখে ধরি $(!A \land !B) \in \Gamma$। তাহলে
$\Gamma \Proves !A \land !B$।
\tagrefs{prfSC/{fol:seq:prv:prop:provability-land-left},prfND/{fol:ntd:prv:prop:provability-land-left}}
অনুসারে $\Gamma \Proves !A$ এবং $\Gamma \Proves !B$।
\olref{prop:mcs-prov-in} অনুসারে $!A \in \Gamma$ এবং $!B \in
\Gamma$, যা প্রয়োজন ছিল।

উল্টো অভিমুখে ধরি $!A \in \Gamma$ এবং $!B \in
\Gamma$। তাহলে $\Gamma \Proves !A$ এবং $\Gamma \Proves !B$।
\tagrefs{prfSC/{fol:seq:prv:prop:provability-land-right},prfND/{fol:ntd:prv:prop:provability-land-right}}
অনুসারে $\Gamma \Proves !A \land !B$। \olref{prop:mcs-prov-in}
অনুসারে $(!A \land !B) \in \Gamma$।}}

\tagitem{defOr}{}{%
\iftag{probOr}{অনুশীলনী।}{%
$(!A \lor !B) \in \Gamma$ যদি এবং কেবল যদি $!A \in \Gamma$ অথবা
$!B \in \Gamma$।

সামনের অভিমুখের বিপরীত-প্রতিজ্ঞাপনের জন্য ধরি, $!A
\notin \Gamma$ এবং $!B \notin \Gamma$। দেখাতে চাই $(!A \lor
!B) \notin \Gamma$। $\Gamma$ সর্বাধিক সঙ্গত বলে $\Gamma
\cup \{!A\} \Proves \lfalse$ এবং $\Gamma \cup \{!B\} \Proves \lfalse$।
\tagrefs{prfSC/{fol:seq:prv:prop:provability-lor-left},prfND/{fol:ntd:prv:prop:provability-lor-left}}
অনুসারে $\Gamma \cup \{(!A \lor !B)\}$ অসঙ্গত। অতএব প্রয়োজনমতো
$(!A \lor !B) \notin \Gamma$।

উল্টো অভিমুখে ধরি $!A \in \Gamma$ অথবা $!B \in
\Gamma$। তাহলে $\Gamma \Proves !A$ অথবা $\Gamma \Proves !B$।
\tagrefs{prfSC/{fol:seq:prv:prop:provability-lor-right},prfND/{fol:ntd:prv:prop:provability-lor-right}}
অনুসারে $\Gamma \Proves !A \lor !B$। \olref{prop:mcs-prov-in}
অনুসারে $(!A \lor !B) \in \Gamma$, যা প্রয়োজন ছিল।}}

\tagitem{defIf}{}{%
\iftag{probIf}{অনুশীলনী।}{%
$(!A \lif !B) \in \Gamma$ যদি এবং কেবল যদি $!A \notin \Gamma$
অথবা $!B \in \Gamma$:

সামনের অভিমুখে ধরি $(!A \lif !B) \in \Gamma$; আর বিপরীতভাবে
ধরি $!A \in \Gamma$ এবং $!B \notin \Gamma$। এই অনুমানগুলি থেকে
$\Gamma \Proves !A \lif !B$ এবং $\Gamma \Proves !A$।
\tagrefs{prfSC/{fol:seq:prv:prop:provability-mp},prfND/{fol:ntd:prv:prop:provability-mp}}
অনুসারে $\Gamma \Proves !B$। কিন্তু \olref{prop:mcs-prov-in}
অনুসারে $!B \in \Gamma$, যা $!B \notin \Gamma$-র বিরোধী।

উল্টো অভিমুখে প্রথমে $!A \notin \Gamma$ ক্ষেত্রটি বিবেচনা করি।
\olref{prop:mcs-either-or} অনুসারে $\lnot !A \in \Gamma$ এবং
তাই $\Gamma \Proves \lnot !A$।
\tagrefs{prfSC/{fol:seq:prv:prop:provability-lif},prfND/{fol:ntd:prv:prop:provability-lif}}
অনুসারে $\Gamma \Proves !A \lif !B$। আবার \olref{prop:mcs-prov-in}
অনুসারে প্রয়োজনমতো $(!A \lif !B) \in \Gamma$ পাই।

এবার $!B \in \Gamma$ ক্ষেত্রটি বিবেচনা করি। তাহলে $\Gamma \Proves !B$
এবং
\tagrefs{prfSC/{fol:seq:prv:prop:provability-lif},prfND/{fol:ntd:prv:prop:provability-lif}}
অনুসারে $\Gamma \Proves !A \lif !B$। \olref{prop:mcs-prov-in}
অনুসারে $(!A \lif !B) \in \Gamma$।}}
\end{enumerate}
\end{proof}

\begin{prob}
\olref[fol][com][mcs]{prop:mcs}-র প্রমাণ সম্পূর্ণ করো।
\end{prob}

\end{document}
